\documentclass[10pt,a4paper]{article}
\usepackage{amsmath,amssymb,amsthm,graphicx}
% Card-local minimal preamble for S7 cards cut from sources/geometry-fl03-hausdorff-dimension-of-limit-sets/pinned-original-source.tex
% (doi:10.5802/jep.241, JEP 10 (2023), CC BY 4.0). Author text untouched.
% Omitted and why: the publisher cedram class (unavailable in the toolchain; the
% assembler supplies article) and the publisher title-page machinery
% (\datereceived/\dateaccepted/\dateepreuves/\datepublished and the smf@ title
% macros retained in the source preamble), per the unavailable-class precedent.
\usepackage{amsmath,amssymb,amsthm}
\usepackage[scr=rsfs,cal=euler]{mathalfa}
\usepackage{xfrac}
\multlinegap0pt
\newcommand{\quand}{\quad\text{and}\quad}
\newcommand{\irr}{\mathrm{irr}}
\newcommand{\red}{\mathrm{red}}
\newcommand{\sym}{\mathrm{sym}}
\newenvironment{smallpmatrix}{\big(\begin{smallmatrix}}{\end{smallmatrix}\big)}
\newcommand\mto{\mathchoice{\longmapsto}{\mapsto}{\mapsto}{\mapsto}}
\newcommand{\psfrac}[2]{\sfrac{(#1)}{#2}}
\newcommand{\spfrac}[2]{\sfrac{#1}{(#2)}}
\newcommand{\Changel}{\mathcode`l="7160}
\newcommand{\Changelback}{\mathcode`l="716C}
\newcommand{\NoChangel}[1]{%
\expandafter\let\csname old\string#1\endcsname=#1
\let#1=\relax
\newcommand{#1}{\mathcode`l="716C\csname old\string#1\endcsname\mathcode`l="7160 }%
}
\NoChangel{\log}
\NoChangel{\ln}
\NoChangel{\lim}
\NoChangel{\limsup}
\NoChangel{\liminf}
\NoChangel{\varlimsup}
\NoChangel{\varliminf}
\NoChangel{\varinjlim}
\NoChangel{\varprojlim}
\usepackage{pgf,tikz}
\usetikzlibrary{arrows}
\newtheorem{theo}{Theorem}[section]
\newtheorem{theorem}[theo]{Theorem}
\newtheorem{coro}[theo]{Corollary}
\newtheorem{lemma}[theo]{Lemma}
\newtheorem{prop}[theo]{Proposition}
\newtheorem{conj}[theo]{Conjecture}
\theoremstyle{definition}
\newtheorem{defi}[theo]{Definition}
\newtheorem{question}[theo]{Question}
\newtheorem{example}[theo]{Example}
\newtheorem{rmk}[theo]{Remark}
\newtheorem*{rem}{Remark}
\newcommand{\proj}{\mathbf{P}}
\newcommand{\R}{\mathbb{R}}
\newcommand{\C}{\mathbb{C}}
\newcommand{\N}{\mathbb{N}}
\newcommand{\Hyp}{\mathbb{H}}
\newcommand{\cC}{\mathcal{C}}
\newcommand{\Cone}{\mathrm{Cone}}
\newcommand{\diag}{\mathrm{diag}}
\newcommand{\SL}{\mathrm{SL}}
\newcommand{\SO}{\mathrm{SO}}
\newcommand{\fa}{\mathfrak{a}}
\newcommand{\frakg}{\mathfrak{g}}
\newcommand{\fp}{\mathfrak{p}}
\newcommand{\fk}{\mathfrak{k}}
\newcommand{\norm}[1]{\Vert #1 \Vert}
\newcommand{\scal}[2]{\langle #1\, |\, #2 \rangle }
\DeclareMathOperator{\Card}{Card}
\DeclareMathOperator{\DimH}{DimH}
\DeclareMathOperator{\Sym}{Sym}
\DeclareMathOperator{\Span}{Span}
\DeclareMathOperator{\ad}{ad}
\let\epsilon\varepsilon
\let\op\iota
\let\phi\varphi
\let\p\varphi
\let\G\Gamma
\let\g\gamma
\let\Log\log
\usepackage{cleveref}
\crefname{originalreference}{original reference}{original references}
\Crefname{originalreference}{Original reference}{Original references}
\makeatletter
\newlabel{coro - ComparisonMuLambdaProximalElement}{{2.36}{source}{}{sourcepointer-1}{}}
\newlabel{coro - ComparisonMuLambdaProximalElement@cref}{{[originalreference][0][]2.36}{[1][0][]source}{}{}{}}
\makeatother
\crefname{originalreference}{original reference}{original references}
\Crefname{originalreference}{Original reference}{Original references}
\makeatletter
\newlabel{lem - ProximalConjugacyClass}{{2.42}{source}{}{sourcepointer-2}{}}
\newlabel{lem - ProximalConjugacyClass@cref}{{[originalreference][0][]2.42}{[1][0][]source}{}{}{}}
\makeatother
\crefname{originalreference}{original reference}{original references}
\Crefname{originalreference}{Original reference}{Original references}
\makeatletter
\newlabel{lem-AMS}{{2.37}{source}{}{sourcepointer-3}{}}
\newlabel{lem-AMS@cref}{{[originalreference][0][]2.37}{[1][0][]source}{}{}{}}
\makeatother
\crefname{originalreference}{original reference}{original references}
\Crefname{originalreference}{Original reference}{Original references}
\makeatletter
\newlabel{lem-AMS guichard wienhard}{{2.38}{source}{}{sourcepointer-4}{}}
\newlabel{lem-AMS guichard wienhard@cref}{{[originalreference][0][]2.38}{[1][0][]source}{}{}{}}
\makeatother
\crefname{originalreference}{original reference}{original references}
\Crefname{originalreference}{Original reference}{Original references}
\makeatletter
\newlabel{prop - Benoist cartan projection continuous}{{2.4}{source}{}{sourcepointer-5}{}}
\newlabel{prop - Benoist cartan projection continuous@cref}{{[originalreference][0][]2.4}{[1][0][]source}{}{}{}}
\makeatother
\title{Entropy and critical exponent agree on the convex hull of the weights}
\author{Olivier Glorieux, Daniel Monclair, Nicolas Tholozan}
\date{doi:10.5802/jep.241; CC BY 4.0}
\begin{document}
\maketitle
\setcounter{section}{2}
\setcounter{theo}{30}
\noindent\textit{Editorial note.} All delivered mathematical text is the author's own, copied from the pinned journal source file sources/geometry-fl03-hausdorff-dimension-of-limit-sets/pinned-original-source.tex (Hausdorff dimension of limit sets for projective Anosov representations); no statement or proof step was rewritten, reordered or removed.\par\medskip
\noindent\textit{Editorial note.} The counted proof is the author's complete argument: All delivered mathematical text is the author own, copied verbatim from the pinned journal source. Provenance note: the candidate comes from the lane-local candidate table lane-s12/candidates-geometry-fl03-hausdorff-dimension-of-limit-sets.json (patch 35). Statement source lines 638-649 with the author own proof at 836-865, which the author places remotely under a proof environment naming this statement (the face was flagged as a label mismatch precisely because its proof is not adjacent). Printed number 2.31 from the source walk, cross-checked against the published anchor of this DOI (theorem 3.1 at statement 1100-1103, whose walk value 3.1 matches the article); verified against the host aux.\par\medskip
\noindent\textit{Editorial note.} The label theo-GeneralComparisonEntropyCriticalExponent names the counted statement, and the proof environment immediately following it in the source is delivered verbatim as its proof.\par\medskip
\noindent\textit{Editorial note.} The counted statement is carried by the author's own theo environment in the article source.\par\medskip
\noindent\textit{Editorial note.} Journal-unit assembly: card-local preamble (author theorem environments and macros verbatim), article class, bibliography backend none; any literature citation carried by the delivered span is disclosed in the notes of this card.\par\medskip
\noindent\textit{Editorial note.} Unresolved citation keys not present in the author's own bibliography: CK, abels1995semigroups, guichard2012anosov.\par\medskip
\noindent\textit{Editorial note.} The bibliographic strings were transliterated to ASCII for the card face (the journal line carries accents and typographic punctuation); the SOURCE.json metadata keeps the original form. No mathematical text is affected.\par\medskip
\noindent\textit{Editorial note.} PRESCREEN RE-READ: Cite-flag re-read: the quality screen takes its citation list from the census row, which spans the whole surrounding section, but the counted statement and the remote proof span contain no citation at all (checked line by line in the pinned source). The card therefore carries no unresolved citation.\par\medskip
\noindent\textit{Editorial note.} Assembly shape note: the card-local preamble is the unit minimal preamble used by the other cards of this source (author package set, author macros, author theorem declarations with the shared counter). The verbatim retained preamble of the unit was rejected because it carries the publisher cedram class lines (datereceived / dateaccepted / dateepreuves and the smf title-page macros); the class is not in the toolchain, and the assembler supplies the article class instead, so the first host assemble stopped on an undefined control sequence at the datereceived line. The counters are seeded so the counted statement prints the number the article gives it.\par\medskip
\noindent\textit{Editorial scope.} the counted statement and its own proof environment only\par\medskip
\noindent\textit{Allowed standard prerequisites (editorial).}\begin{itemize}
\item The definitions of the entropy and of the critical exponent of a Th.... Supplies background that the delivered proof uses. The entry is declared as a dependency; it is not re-proved on the card, and the load-bearing steps of the counted statement are carried out inside the delivered span.
\end{itemize}\medskip
\ifdefined\hypertarget\hypertarget{sourcepointer-1}{}\fi
\noindent\textit{Source reference.} Original reference 2.36: sources/geometry-fl03-hausdorff-dimension-of-limit-sets/pinned-original-source.tex, source lines 692--692. Navigation pointer invoked at line(s) 847: the counted text refers to this environment of the same journal unit, which is not retained in this card, so the reference is declared as a pointer carrying the number the article itself prints for it (the section and the place of the statement inside that section, derived from the shared section-scoped theorem counter of the source); the cited item is an allowed input and none of it is reproduced inside the span.\par\medskip
\ifdefined\hypertarget\hypertarget{sourcepointer-2}{}\fi
\noindent\textit{Source reference.} Original reference 2.42: sources/geometry-fl03-hausdorff-dimension-of-limit-sets/pinned-original-source.tex, source lines 817--817. Navigation pointer invoked at line(s) 847: the counted text refers to this environment of the same journal unit, which is not retained in this card, so the reference is declared as a pointer carrying the number the article itself prints for it (the section and the place of the statement inside that section, derived from the shared section-scoped theorem counter of the source); the cited item is an allowed input and none of it is reproduced inside the span.\par\medskip
\ifdefined\hypertarget\hypertarget{sourcepointer-3}{}\fi
\noindent\textit{Source reference.} Original reference 2.37: sources/geometry-fl03-hausdorff-dimension-of-limit-sets/pinned-original-source.tex, source lines 736--736. Navigation pointer invoked at line(s) 855: the counted text refers to this environment of the same journal unit, which is not retained in this card, so the reference is declared as a pointer carrying the number the article itself prints for it (the section and the place of the statement inside that section, derived from the shared section-scoped theorem counter of the source); the cited item is an allowed input and none of it is reproduced inside the span.\par\medskip
\ifdefined\hypertarget\hypertarget{sourcepointer-4}{}\fi
\noindent\textit{Source reference.} Original reference 2.38: sources/geometry-fl03-hausdorff-dimension-of-limit-sets/pinned-original-source.tex, source lines 742--742. Navigation pointer invoked at line(s) 864: the counted text refers to this environment of the same journal unit, which is not retained in this card, so the reference is declared as a pointer carrying the number the article itself prints for it (the section and the place of the statement inside that section, derived from the shared section-scoped theorem counter of the source); the cited item is an allowed input and none of it is reproduced inside the span.\par\medskip
\ifdefined\hypertarget\hypertarget{sourcepointer-5}{}\fi
\noindent\textit{Source reference.} Original reference 2.4: sources/geometry-fl03-hausdorff-dimension-of-limit-sets/pinned-original-source.tex, source lines 393--393. Navigation pointer invoked at line(s) 855: the counted text refers to this environment of the same journal unit, which is not retained in this card, so the reference is declared as a pointer carrying the number the article itself prints for it (the section and the place of the statement inside that section, derived from the shared section-scoped theorem counter of the source); the cited item is an allowed input and none of it is reproduced inside the span.\par\medskip
\begin{theo} \label{theo-GeneralComparisonEntropyCriticalExponent}
Let $\Gamma$ be a $\Theta$-Anosov subgroup of $G$. Then
\begin{itemize}
\item $h_\phi(\Gamma) \leq \delta_\phi(\Gamma)$ for all $\phi\in D^*(\Theta)\cap \fa_+^*(\Theta)$,
\item $h_\phi(\Gamma) = \delta_\phi(\Gamma)$ for all $\phi \in W(\Theta)\cap \fa_+^*(\Theta)$.
\end{itemize}

If $\G$ is moreover Zariski dense in $G$, then
\begin{itemize}
\item $h_\phi(\G) =\delta_\phi(\G)$ for all $\phi \in D^*(\Theta)\cap \fa_+^*(\Theta)$.
\end{itemize}
\end{theo}

\begin{proof}[Proof of Theorem \ref{theo-GeneralComparisonEntropyCriticalExponent}]

\Changel

Let $\varphi$ be an element of $D^*(\Theta)\cap \fa_+^*(\Theta)$. Since $\Gamma$ is $\Theta$-Anosov and $\phi \in \fa_+^*(\Theta)$, there exist $\alpha, A>0$ such that for all $\g \in \G$,
\[\alpha \vert \g \vert - A \leq \phi(\mu(\g)) \leq \frac{1}{\alpha}\,\vert \g \vert + A\quand
\alpha l([\g]) - A \leq \phi(\lambda(\g)) \leq \frac{1}{\alpha}\,l([\g]) + A.\]
We first prove the inequality $h_\varphi \leq \delta_\varphi$, then the reverse inequality.

\Changelback

Fix $r>0$ given by Corollary \ref{lem - ProximalConjugacyClass}, then $\epsilon>0$ given by Corollary \ref{coro - ComparisonMuLambdaProximalElement}, then $D$ as in Corollary \ref{lem - ProximalConjugacyClass}. With these choices, every conjugacy class $[\gamma]$ with $\phi(\lambda(\g)) \geq \psfrac{D-A}{\alpha}$ contains at least one $(\Theta,r,\epsilon)$-proximal element.

Define
\[\Gamma_\phi(R) =\{\g \in \G \mid \phi(\mu(\gamma)) \leq R\}.\]
For every $R>0$ and every conjugacy class $[\g] \in [\G]$ of infinite order such that $\psfrac{D-A}{\alpha}\leq \varphi(\lambda(\g)) \leq R$, we can find $\g' \in [\g]$ which is $(\Theta,r,\epsilon)$-proximal by Corollary~\ref{lem - ProximalConjugacyClass}. By Corollary \ref{coro - ComparisonMuLambdaProximalElement}, since $\varphi \in D^*(\Theta)$, there exists a constant $C$ (independent of $R$) such that $\gamma' \in \Gamma_\phi(R+C)$. We deduce that
\[\Card \Gamma_\phi(R+C) \geq \Card \{[\g] \in [\Gamma] \textrm{ of infinite order } \mid \psfrac{D-A}{\alpha}\leq \varphi(\lambda(\g))\leq R\}.\]
The inequality $\delta_\varphi \geq h_\varphi$ easily follows.

To show $\delta_\varphi \leq h_\varphi$, assume first that $\G$ is Zariski dense. Take $r>0$ given by Abels--Margulis--Soifer's Lemma \ref{lem-AMS}, then $\epsilon$ such that Lemma \ref{coro - ComparisonMuLambdaProximalElement} applies to $(r,\epsilon)$\nobreakdash-loxo\-dromic elements. Lemma \ref{lem-AMS} then gives us a finite subset $F\subset \G $ such that for any $\g\in \G$, there exists $f\in F$ such that $f\g$ is $(r,\epsilon)$-loxodromic. By Proposition~\ref{prop - Benoist cartan projection continuous} and Corollary \ref{coro - ComparisonMuLambdaProximalElement}, there is a constant $C'$ depending on $F, r$ and $\epsilon$ such that \hbox{$|\varphi(\lambda(f\g)) - \varphi(\mu(\g))|\leq C'$}. We thus have
\[
\Card\Gamma_\phi(R) \leq \Card(F) \cdot \Card\{ \g \in \Gamma \mid \phi(\lambda(\g)) \leq R+C', \textrm{ $\g$ is $\epsilon$-loxodromic}\}.\]
Now, by Corollary \ref{lem - ProximalConjugacyClass}, every conjugacy class $[\g]$ such that $\phi(\lambda(\gamma)) \leq R+C'$ contains at most $\sfrac{C(R+A+C')}{\alpha}$ elements that are $(r,\epsilon)$-loxodromic, for the constant $C$ given by Corollary \ref{lem - ProximalConjugacyClass}.

We thus obtain
\[\Card \Gamma_\phi(R) \leq \frac{C(R+A+C')}{\alpha} \Card F \cdot \Card \{[\g]\in [\Gamma] \mid \varphi(\lambda(\g))\leq R+C\},\]
from which the inequality $\delta_\varphi \leq h_\varphi$ easily follows.

If $\Gamma$ is not assumed Zariski dense, we can still apply Lemma \ref{lem-AMS guichard wienhard} and choose $r>0$ for which every element $\g\in \G$ can be made $(\Theta,r,\epsilon)$-proximal by multiplying it by an element $f$ in a finite set. By Corollary \ref{coro - ComparisonMuLambdaProximalElement}, we do have $|\phi(\lambda(f \g)) - \phi(\mu(\g))| \leq C'$ provided that $\phi$ belongs to $W(\Theta)$. The rest of the proof works in a similar way and we eventually obtain the inequality $\delta_\phi \leq h_\phi$ for $\phi \in W(\Theta)$.
\end{proof}

\end{document}
